%!TEX root = ../main.tex
\begin{abstractko}
리그 오브 레전드(LoL)에서 교전(한타)은 경기의 흐름을 결정하는 핵심 사건이지만, 기존 승패 예측 연구는 경기 전체를 예측 단위로 삼거나 게임 클라이언트에서만 얻을 수 있는 초 단위 리플레이 관측을 전제하였다. 본 논문은 누구나 접근 가능한 Riot Match-V5 공개 텔레메트리, 즉 60 초 간격의 상태 프레임과 밀리초 단위의 사건 기록만으로 교전 결과를 어디까지 예측할 수 있는지를 체계적으로 규명한다.

첫째, 킬 사건만을 씨앗으로 하는 결정론적 교전 국소화 알고리즘을 제안한다. 시간 군집화(18 초), 완전 연결 기준의 공간 지름 분할(4{,}000 단위), 온셋 검증(팀당 생존 2 명 이상, 1{,}800 단위), 인접 후보 병합(15 초 / 2{,}000 단위)의 네 단계로 \numMatches{} 경기의 한국 서버 상위 티어 랭크 경기에서 약 100 만 건의 교전을 사람의 주석 없이 구축하였다. 둘째, 킬·셧다운·오브젝트·현상금의 도메인 가치를 softmax 주의로 가중하는 교환가치 라벨을 정의하였다. 셋째, 플레이어 노드 76 차원, 전역 26 차원, 사건 44 차원의 특징을 6 개의 5 초 구간으로 구성하되 60 초 프레임을 계단 유지하고 온셋 이후 정보를 결코 사용하지 않는 무누출 계약을 명시하였다. 넷째, LightGBM, 정합 입력 다층 퍼셉트론, 순환·주의·합성곱·상태공간 시퀀스 모델, 그래프 신경망, 시공간 그래프 신경망, 계층적 융합 모델 등 25 종 이상의 모델을 동일한 시간순 패치 홀드아웃(학습 15.14 / 검증 15.15 / 시험 15.16)과 세 시드, 부트스트랩 신뢰구간, DeLong 검정, Holm--Bonferroni 보정 아래에서 비교하고, 일곱 가지 도메인 지식 처치를 5 단계 절제 프로토콜로 평가하였다.

공학적 테이블 요약 위의 LightGBM이 시험 AUC $\aucLGBM{}$로 가장 우수하였고, 같은 입력의 MLP는 $\aucMLP{}$, 시퀀스·그래프·사건 기반 신경망 표현은 $.569$--$\aucBiGRU{}$에 머물렀다. 정합 입력 비교는 이 격차의 절반이 학습기 계열에서, 나머지 절반이 입력 표현에서 비롯됨을 보였다. 예측 가능성은 게임 국면(초반 $\aucEarly{}$, 중반 $\aucMid{}$, 후반 $\aucLate{}$)과 온셋 시점 골드 격차(대등 $\aucClose{}$, 보통 $\aucModerate{}$, 일방 $\aucOneSided{}$)에 따라 단조 증가하였으며, 대등한 교전에서는 우연 수준에 근접하였다. 이는 분 단위 공개 텔레메트리의 정보 입도가 교전 결과 예측에 부과하는 상한을 시사한다. 본 논문은 검출기의 사람 주석 대비 타당성 검증 프로토콜과 코드--논문 감사 결과를 함께 제시하며, 코드는 MIT 라이선스로 공개한다.
\end{abstractko}
